%
% ---------- header -----------------------------------------------------------
%
% project       kaneton
%
% license       kaneton
%
% file          /home/mycure/kaneton/view/book/development/development.tex
%
% created       julien quintard   [mon may 14 19:56:45 2007]
% updated       julien quintard   [wed may 16 21:42:55 2007]
%

%
% path
%

\newcommand{\path}{../..}

%
% template
%

%
% ---------- header -----------------------------------------------------------
%
% project       kaneton
%
% license       kaneton
%
% file          /home/mycure/kaneton/view/template/book.tex
%
% created       julien quintard   [wed may 16 18:16:50 2007]
% updated       julien quintard   [sun jun  3 12:00:37 2007]
%

%
% class
%

\documentclass[10pt,a4wide]{book}

%
% packages
%

\usepackage[english]{babel}
\usepackage[T1]{fontenc}
\usepackage{a4wide}
\usepackage{fancyheadings}
\usepackage{multicol}
\usepackage{indentfirst}
\usepackage{graphicx}
\usepackage{color}
\usepackage{xcolor}
\usepackage{verbatim}

\usepackage{aeguill}

\usepackage[Lenny]{\path/package/fncychap}

\pagestyle{fancy}

%
% -
%

\renewcommand{\-}{\vspace{\parskip}}

%
% toc
%

\newcommand{\toc}
  {
    \tableofcontents
    \setlength{\footrulewidth}{0.3pt}
    \setlength{\parindent}{0.3cm}
    \setlength{\parskip}{2ex plus 0.5ex minus 0.2ex}
  }

%
% logos
%

\newcommand{\logos}
  {
    \begin{center}
      \includegraphics[scale=0.8]{\path/logo/kaneton.pdf}
    \end{center}
  }

%
% colors
%

\definecolor{functioncolor}{rgb}{0.40,0.00,0.00}
\definecolor{commandcolor}{rgb}{0.00,0.00,0.40}
\definecolor{verbatimcolor}{rgb}{0.00,0.40,0.00}
\definecolor{noticecolor}{rgb}{0.87,0.84,0.02}

%
% function
%

\newcommand\function[3]{
  \begin{tabular}{p{0.2cm}p{13.8cm}}
  & {\color{functioncolor}\textbf{#1}}#2
  \end{tabular}

  \begin{tabular}{p{1cm}p{13cm}}
  & #3
  \end{tabular}}

%
% align
%

\newcommand\align[1]{
  \\ & \hspace{#1}}

%
% argument
%

\newcommand\argument[1]{\textit{#1}}

%
% command
%

\newcommand\command[2]{
  \begin{tabular}{p{0.2cm}p{13.8cm}}
  & {\color{commandcolor}\textbf{#1}}
  \end{tabular}

  \begin{tabular}{p{1cm}p{13cm}}
  & #2
  \end{tabular}}

%
% notice
%

\newcommand\notice[1]{
  {\color{noticecolor}\textbf{Notice}}

  \begin{tabular}{p{0.2cm}p{13.8cm}}
  & #1
  \end{tabular}}

%
% subsubsubsection
%

\newcommand\subsubsubsection[1]{\textbf{#1}}

%
% example
%

\newcommand\example[1]{
  \textit{Example:}

  \begin{tabular}{p{0.2cm}p{13.8cm}}
  & \textit{#1}
  \end{tabular}}

%
% warning XXX
%

\renewcommand{\familydefault}{\sfdefault}

%
% verbatim stuff
%

\makeatletter

\renewcommand{\verbatim@font}
  {\ttfamily\footnotesize\selectfont}

\def\verbatim@processline{
  \hskip15ex{\color{verbatimcolor}\the\verbatim@line}\par
}

\makeatother

%
% header
%

\rhead{}
\rfoot{\scriptsize{The kaneton microkernel project}}

\date{\scriptsize{\today}}


%
% header
%

\lhead{\scriptsize{The kaneton microkernel project reference}}
\rhead{}

%
% title
%

\title{The kaneton microkernel reference $_{beta}$
       \logos}

%
% authors
%

\author{\small{Julien Quintard}}

%
% document
%

\begin{document}

%
% title
%

\maketitle

%
% --------- text --------------------------------------------------------------
%

This document describes how people can contribute to the kaneton microkernel
project and the rules they have to follow.

This document must be read by everyone involved in the kaneton microkernel
project.

All the kaneton documents are available on
the official website
  \footnote{http://www.kaneton.org}.

%
% toc
%

\tableofcontents

%
% chapters
%

%
% ---------- header -----------------------------------------------------------
%
% project       kaneton
%
% license       kaneton
%
% file          /home/mycure/kaneton/view/book/development/introduction.tex
%
% created       julien quintard   [thu may 17 12:46:30 2007]
% updated       julien quintard   [fri jun 15 11:48:28 2007]
%

%
% ---------- introduction -----------------------------------------------------
%

\chapter{Introduction}
\label{chapter:environment}

In this chapter, the kaneton microkernel project is briefly introduced
in order to emphasize some of its characteristics that makes contribution
specific in several ways.

\newpage

%
% ---------- text -------------------------------------------------------------
%

kaneton is a educational purpose microkernel project. This project aims
at providing a very clear, commented and maintainable microkernel source
code in order to allow people interested in operating systems internals
to look at the source code and understand it very quickly.

The kaneton project is basically composed of the source code of the
microkernel itself, scripts to perform complex tasks and various documents
from design papers to lecture materials.

The most important thing to remember is that the whole project is intended
to be understood as well as possibly maintained by everyone. As a result,
contributions must comply with the level of clarity expected by the project.

These rules are discussed in this paper in order to inform every new
contributor of what makes a good contribution.

The remaining of this document is organised as follows. \textit{Chapter
\ref{chapter:history}} draws the history of the kaneton microkernel project.
\textit{Chapter \ref{chapter:source tree}} introduced the kaneton project
organisation through the source code hierarchy. Next, \textit{Chapter
\ref{chapter:community}} describes how a contributor should behave in a
development community. \textit{Chapter \ref{chapter:rules}} introduces the
general rules which apply to any context around the kaneton project. Then,
the tools inherent to the kaneton project are listed in \textit{Chapter
\ref{chapter:tools}} with some guidelines about how to use them properly.
\textit{Chapter \ref{chapter:languages}} explicitly describes languages rules
informing the developer of the coding style to respect. \textit{Chapter
\ref{chapter:people}} draws a list of the people in charge for the different
parts and tools of the project. Finally, \textit{Chapter
\ref{chapter:licenses}} contains information about the licenses related to
the kaneton microkernel project.

%%
% ---------- header -----------------------------------------------------------
%
% project       kaneton
%
% license       kaneton
%
% file          /home/mycure/kaneton/view/book/development/community.tex
%
% created       julien quintard   [sun may 20 18:08:17 2007]
% updated       julien quintard   [tue may 20 17:03:59 2008]
%

%
% ---------- community --------------------------------------------------------
%

\chapter{Community}
\label{chapter:community}

This chapter discusses what is a community and how contributors must integrate
the kaneton community.

\newpage

%
% ---------- text -------------------------------------------------------------
%

kaneton can obviously be considered as an open source community although the
produced soure code is actually not open source.

Driving an open source community is complicated since people have different
personal goals at working on a free project. Some people contribute for
the knowledge, other for building the next generation system, other to
provide free open source softwares, other to become famous \etc{}

kaneton is a community driven microkernel that acts with the best interest of
the students at heart. Rules and regulations that keep the project moving
forward are fundamental even if the size of the kaneton community is
relatively small, for now.

Indeed, the main objective of the kaneton project remains to be as
understandable as possible in order to lead students to implement parts of
it very quickly.

The remaining of this chapter draws a list of rules contributors must agree
to respect.

% objective

\subsubsection{Objective}

kaneton aims at providing a powerful, understandable and maintainable
microkernel. This objective must be kept in mind of every contributor
since many design and implementation were/are/will be made according to
this precise objective.

Note that the kaneton microkernel does not intend to be a desktop operating
system nor an as optimised as the Linux operating system. Every contributor
should be well-aware of that in order to avoid behaviours stating that a
feature is fundamental or useless for performance concerns, for instance.

This rule does not prohibit people to suggest ideas but instead regulates
behaviours of people who wants to change major design and/or implementation
choices for bad reasons.

% behaviour

\subsubsection{Behaviour}

Open source projects does not mean constraint-free projects. The kaneton
people, whilst being relatively young, try to act for the project's good
by behaving remarkably in the kaneton community.

Therefore, contributors are asked to do the same by avoiding some bad/young
behaviours.

\begin{enumerate}
  \item
    Follow the rules. People who do not respect these rules could be banned
    from the kaneton project.
  \item
    Avoid the \textit{cowboy} behaviour consisting for a contributor to
    implement a feature without discussing about its usefulness with the
    community first. Another effect of this behaviour can be to distract
    the contributors from its major focuses.
  \item
    Always act and think in the project interest rather than your personal
    interest.
  \item
    Respect the other kaneton people, especially the ones who have worked
    on this project for a long time and who made this whole project possible.
    When people disagree, they are asked to do it respectfully.
  \item
    Take your responsability when you realise that you did something wrong:
    insults, mistakes in an implemented feature \etc{}
  \item
    \textit{``The Perfect is the Enemy of the Good!''}: even nice
    contributors can unintentionally do bad things by being perfectionists
    and/or too much into the project and/or obsessed with process.
  \item
    ... \textit{Politness}, \textit{Respect}, \textit{Trust} and
    \textit{Humility} are the key qualities that make a good contributor in
    any community.
\end{enumerate}

% communication

\subsubsection{Communication}

The communication mainly takes two forms in the kaneton microkernel project:
the \name{mailing-list} for internal communication and the \name{kaneton
website} for external communication. The \name{Developers Wiki} is another
source of communication as well as the commit logs \etc{}

The rules related to these tools are described in \reference{Chapter
\ref{chapter:tools}} and will therefore not be discussed here.

Every contributor must take the time to communicate in the mailing-list as well
as through the public website. Indeed, kaneton people must, frequently,
briefly describe what they are working on in order to inform the other
contributors who are not aware of everyone's current work. Note that these
kind of messages are very different from messages generated by repository
commits. Indeed, while these commit messages indicate a modification, they
do not describe the whole work behind them.

Additionally, contributors can communicate informing the kaneton community of
their unavailability for the next two months, for instance. This behaviour
allows people to be aware that some tasks will not be done because the
contributors in charge of it cannot work at this moment.

Althoug people are highly welcomed to communicate, some rules apply to
avoid further problems.

First, any new contributor should obviously read the kaneton documents and
then go through the mailing-list archives before asking anything which has
already been discussed, unless the developer knows exactly what he/she is
talking about. Indeed, asking too many questions about the source code is a
form of disrespect to the other contribors. Moreover, many things can be
found out just looking at the kaneton documents and/or source code.

Although people are asked to communicate, people are also asked to act
respectully. Contributors should not respond to every message in every
discussion, this is a ridiculous behaviour. Instead, every developer should
carefully read the discussion, think about its response and then write a clear
message stating his point of view, ideas \etc{}

Depending on the contributor status, reading the mailing-list frequently is
absolutely fundamental as some people rely on other contributors' decisions,
advices \etc{}

Finally, the mailing-list must be considered as the official internal
communication medium. If contributors previously had a private conversation,
in real-life or on \name{IRC} for instance, the discussion must be reported
on the mailing-list so that everyone can take these new ideas into account.

% work

\subsubsection{Work}

Working on the kaneton microkernel project does not imply low-level programming
all the time. Indeed, the kaneton project is composed of two parts: the
kaneton microkernel research project and the kaneton educational project.

Although the microkernel research project requires highly skilled programmers,
it also needs documentations and some tools for performing important tasks as
diverse as generating the prototypes, testing the microkernel behaviour,
generating the documentation \etc{}

The educational project essentially needs documentation, lecture materials
and tools for managing the project: testing the students' implementation,
checking if some students cheated and many others.

This means that kaneton people must contribute to every type of task
that need to be done. Also, contributors are asked to well document
any work they have done including source code comments but also through
kaneton official documents which are then made available on the website.

% supervisor

\subsubsection{Supervisor}

A supervisor is attached to each new contributor for a certain, not fixed,
period of time. The role of this supervisor is to advise, correct and
encourage the newcomer so that it integrates well the kaneton community.

The supervisor will be someone having a well understanding of the new
contributor's project. A high level of communication is expected between
the supervisor and the contributor. A direct phone communication is highly
recommanded through \name{Skype} for instance.

% trust

\subsubsection{Trust}

A new contributor joining the kaneton project must acquire the trust of
the community. Therefore, the contributor first does not get any access
to restricted tools and must submit patches to its supervisor who review
them, taking care of advising the newcomer of its mistakes.

At the end of the test period, the community decides whether the contributor
is helpful to the project or not. Then, the contributor is either granted
of full access to kaneton tools or eliminated from the project.

%%
% ---------- header -----------------------------------------------------------
%
% project       kaneton
%
% license       kaneton
%
% file          /home/mycure/kaneton/view/book/development/source-tree.tex
%
% created       julien quintard   [thu may 17 22:41:36 2007]
% updated       julien quintard   [sun may 20 18:09:37 2007]
%

%
% ---------- source tree ------------------------------------------------------
%

\chapter{Source Tree}

In this chapter we will briefly describe the kaneton microkernel project
source tree.

\newpage

%
% ---------- text -------------------------------------------------------------
%

The kaneton microkernel reference source tree looks like the following
listing:

\begin{verbatim}
cheat/
environment/
export/
configure/
kaneton/
library/
license/
tool/
test/
transcript/
view/
\end{verbatim}

% cheat/

\subsection*{cheat/}

Since the kaneton microkernel is implemented by students, the kaneton
people need to check whether students are cheating by re-using parts of
previous years projects or other kernel source code available on the
\textit{Internet}.

To avoid cheating, kaneton people developed a software checking for
commonalities between different source codes.

The \textit{cheat/} kaneton subdirectory contains everything necessary
to perform such tasks including the software, the different source codes
found on the internet, the kaneton students' implementation from the
previous years etc.

% environment/

\subsection*{environment/}

This directory contains everything necessary to the kaneton development
environment.

The kaneton development environment allows different developers to
interact on the development of the same microkernel in a pretty easy way.

The development environment aims at providing developers to possibility to
work in a collaborative manner without interfering with each other. These
developers are likely to run different operating systems on different
microprocessors. In addition, the kaneton microkernel can be targeted for
different microprocessor architectures. The development environment was
introduced to cope with these combinaisons by providing profiles, each
profile describing the behaviour of a component: underlying operating system,
target architecture, user-specific stuff etc.

As a result, each developer can use a different operating system and
microprocessor architecture with his own specific compiling flags, kaneton
parameters etc. without modifying another developer's configuration.

% export/

\subsection*{export/}

The \textit{export/} directory contains scripts used to generate a kaneton
tarball in order to be distributed to the students at the beginning of the
kaneton project.

Indeed, these scripts rearrange the kaneton hierarchy hidding some
important directories the students do not need to be aware of.

Moreover some source code parts are removed since the students then have to
rewrite these piece of code.

% configure/

\subsection*{configure/}

This directory contains everything necessary for configuring its own
kaneton microkernel development environment through the compiling process
to the boot system.

Any new contributor should first look at this directory.

% kaneton/

\subsection*{kaneton/}

This directory is the most important of the project since it contains
the whole microkernel source code.

The directory is composed of three important subdirectories: \textit{core/},
\textit{platform} and \textit{architecture}. These subdirectories are described
next.

% kaneton/core/

\subsection*{kaneton/core/}

This directory contains the kaneton core source code.

The directory is divided as shown below:

\begin{verbatim}
as/
region/
sched/
segment/
set/
task/
thread/
[...]
\end{verbatim}

Each directory represents a kaneton core manager. For more information on
the kaneton core, please refer to the appropriate document:
\textit{The kaneton microkernel :: reference implementation}

% kaneton/platform/

\subsection*{kaneton/platform/}

This directory contains everything in relation with what the kaneton
microkernel project calls a \textbf{platform}. The platform represents the
board supporting the devices: microprocessor, memory, peripherals etc.

This directory obviously contains subdirectories for each platform
supported by the kaneton microkernel.

% kaneton/architecture

\subsection*{kaneton/architecture}

The \textit{architecture} directory contains the source-code related to
the microprocessor architecture.

This directory is composed of subdirectories, each one representing an
architecture: \textit{ia32}, \textit{mips64} etc.

% library/

\subsection*{library/}

This directory contains the libraries used either by the kaneton microkernel
itself or by the kaneton microkernel servers. This directory especially
contrains the standard \textit{kaneton C library}.

% license/

\subsection*{license/}

This directory contains the licenses used for any program or document
in relation with the kaneton microkernel project.

Each student has to read and accept the kaneton license before implementing
or even using the kaneton microkernel project.

% tool/

\subsection*{tool/}

This directory contains additional files used by the kaneton development
environment.

% test/

\subsection*{test/}

Since the kaneton microkernel is used as a material for operating system
courses, the kaneton microkernel reference which is the basis of students'
work is to be extremely reliable.

The kaneton project therefore contains a set of tools in order to validate
the kaneton reference implementation behaviour. These tools are also used
for evaluating the correctness of the students' implementation.

The \textit{test/} directory so contains the set kaneton scripts and tests
for validating a kaneton microkernel implementation.

% transcript/

\subsection*{transcript/}

This directory contains real-time recorded sessions. These sessions can be
replayed showing how to use the kaneton microkernel project and more
specifically the kaneton development environement.

% view/

\subsection*{view/}

This directory contains all the kaneton documents including kaneton
administrative documents, exams documents, lectures materials, kaneton papers
and books etc.

Additionally, scripts are provided in order to very easily build and
display these documents.

%%%
%% licence       kaneton licence
%%
%% project       kaneton
%%
%% file          /home/mycure/kaneton/view/books/kaneton/development-environment.tex
%%
%% created       matthieu bucchianeri   [mon jan 30 17:31:26 2006]
%% updated       julien quintard   [sun jul 30 23:01:51 2006]
%%

%
% development environment
%

\chapter{Development Environment}

In this chapter we will study the kaneton development environment
to understand how it works and how to extend it with new users and
machines profiles.

\newpage

%
% text
%

Over the years, the kaneton microkernel evolved, starting with a
very simple introduction to low-level programming and finally
to a complete microkernel development.

Going always further implies many modifications. Indeed, the students
cannot take time to develop their own development environment because
developing a development environment is a whole project by itself.

We wanted to lead students to a complete microkernel development and then
to introduce the parallel and distributed programming. This would not
be possible if they had to develop an entire development environment.

So kaneton people decided to give the students a complete development
environment.

The kaneton development environment is composed of makefiles, shell
scripts and configuration files and together provide multi-users cooperation,
with many operating systems on many architectures.

A development environment describes binaries, compiling options, linking
options etc.. to use for the kaneton microkernel project.

The whole kaneton development environment needs exactly two
fundamental tools to work. The first one is GNU make, used to build
powerful makefiles, and the second one is GNU bash, used to write
shell scripts.

If an operating system has these two tools, then kaneton can be developed
on it.

%
% profiles
%

\section{Profiles}

The kaneton development environment permit different users, with different
development parameters and configuration files, running different operating
systems on different architectures, to develop together on the same
microkernel project.

To offer these functionalities, the kaneton development environment is
composed of profiles, machine profiles which describe operating systems
behaviours and user profiles which describe users configurations.

From these profiles, the kaneton development environment will generate
development environment files containing the whole kaneton environment
stuff including variables, shell functions, make functions etc..

The three kaneton development environment files are located in the
\textit{env/} directory.

The first one contains static kaneton development environment information
and is called \textit{.env.conf}. This file contains absolute paths
to different kaneton directories and kaneton general information.

The two others called \textit{.env.mk} and \textit{.env.sh} will be generated
by the kaneton development environment system. These files contain
everything necessary to kaneton makefiles and kaneton shell scripts
including variables and functions.

The kaneton development environment system uses shell environment variables
the user has to set before using the system:

\begin{itemize}
  \item
    \textit{\$KANETON\_USER} holds the user profile name.

    This shell environment variable will be used by the kaneton development
    environment system to retrieve the user profile.
  \item
    \textit{\$KANETON\_MACHINE} holds the machine profile name.

    This value will be used to retrieve the machine profile.
  \item
    \textit{\$KANETON\_SHELL} holds the GNU bash binary path.

    This path will be used to launch the very first shell script used to
    generate the current development environment.
\end{itemize}

%
% machines
%

\subsection{Machines}

kaneton people wanted a powerful system to provide users an easy way
to use kaneton on any operating system and any microprocessor architecture.

Indeed, many students and developers use Apple's MacOS running on PowerPC
microprocessors while other use Solaris running on UltraSparc microprocessors
and finally the others generally use Linux or BSD operating systems on
Intel microprocessors.

kaneton people wanted all the developers to contribute to the same
kaneton implementation without difficulties.

To make it possible, the kaneton microkernel had to compile on
every operating systems. To do so, machine profiles were introduced to
describe operating systems behaviour.

A machine profile is precisly a couple operating system, running
microprocessor architecture and kaneton target microprocessor architecture.

The naming scheme used for machine profiles is the operating system name,
a dash, the running on microprocessor, a dot and the target microprocessor
architecture for which kaneton is built.

Note that in the case of a same running on microprocessor and target
microprocessor, only the first naming part is necessary.

For example, a correct machine profile name would be
\textit{macos-ppc.ia32} describing the machine profile composed of the MacOS
operating system running on PowerPC microprocessor architecture and generating
kaneton binary object files for the Intel Architecture 32-bit. Another correct
machine profile name would be \textit{freebsd} describing the behaviour of
the FreeBSD operating system running on any microprocessor and generating
kaneton binary object files for the same microprocessor's architecture.

The machine profiles are located in the \textit{env/machines/} directory.
A subdirectory is used for every machine profiles.

Notice that the kaneton development environment will use the shell
environement variable \textit{\$KANETON\_MACHINE} to locate the
machine profile to use.

Then the current machine profile's files are always located in the
\textit{env/machines/\$KANETON\_MACHINE/} directory.

A machine profile describes the operating system behaviour for very
generic operations like compiling a file, assembling a file, building
a link between two files, removing a file, changing the current directory,
displaying a message etc..

As seen earlier, the kaneton development environment describes two
behaviours, one for the makefiles, and another for the shell scripts.

The files located in each machine profile are listed below. Note that
machine profiles may add any necessary file in their directory.

\begin{verbatim}
clean.sh
critical.sh
init.sh
machine.conf
machine.mk
machine.sh
\end{verbatim}

The files \textit{critical.sh}, \textit{init.sh} and \textit{clean.sh} are
used to install and uninstall the development environment for this
machine profile.

More precisly, the \textit{critical.sh} shell script must generate
the development environment fundamental files \textit{.env.mk} and
\textit{.env.sh} located in the \textit{env/} directory.

Then, the \textit{init.sh} is called to install machine specific stuff
like links to machine dependent directories. This shell script should also
check for the presence of binaries used by the machine profile.

Finally, the \textit{clean.sh} is called to uninstall the machine
specific stuff.

The \textit{machine.conf} file describes the machine profile configuration
including general variables, general flags etc.. for the machine profile.

The file \textit{machine.mk} describes the machine behaviour for
make operations. This file must implement the whole kaneton make interface.

The file \textit{machine.sh} describes the machine behaviour for
shell operations. This file must implement the whole kaneton shell
interface.

Note that it might exist partial behaviours while the user willing use
this machine profile must implement the other part.

Then, with all these files, an operating system can be described to
perform very basic make and shell operations generating binary object
files for the current microprocessor architecture or for another specific
one.

%
% users
%

\subsection{Users}

Each kaneton developer has his own profile located in the kaneton
hierarchy \textit{env/users/\$KANETON\_USER/}.

Note that this shell environment variable must be composed of a
first name, a dot and a last name; using nickname is not allowed.

For example, a correct user name would be \textit{julien.quintard} and
a bad one would be \textit{ChIcHe}.

The user profile directory contains configuration files used to parameter
the developement environment and the kaneton microkernel. Since every
users has his own profile, the current developer will use his own user
profile configuration and will not interfer with other developers.

The files located in each user profile directory are listed below.

\begin{verbatim}
  conf.c
  conf.h
  kaneton.conf
  modules.conf
  user.conf
  user.mk
  user.sh
\end{verbatim}

The file \textit{conf.c} is not used in the current development environment
implementation.

The file \textit{conf.h} parameterizes the kaneton microkernel specifying
for example the algorithm to use in the segment and region managers,
but is also used to activate for example the debug modes etc..

The file \textit{kaneton.conf} is the runtime kaneton microkernel
configuration file. This file describes the kaneton servers to launch
at the operating system startup. Moreover this file is used to parameter
the kaneton microkernel and servers in a dynamic way.

The file \textit{modules.conf} contains the module files to pass to
kaneton. Note that the first module in this file will be considered as
the kaneton bootloader and the second one as the kaneton core.

The file \textit{user.conf} is the most important since it is used
to configure the whole kaneton development environment. This file uses
the GNU make syntax. Via this file, the developer can parameter the compiling
flags, the boot device modes, the display options, the microkernel
files to use etc..

The \textit{user.mk} makefile is used to redefine or complete the machine
makefile behaviour.

Indeed, a developer could want not to use the general machine behaviour but
instead, prefer redefine parts of it. Then the user can override parts of
the machine profile's makefile behaviour using this file to redefine makefile
functions.

The \textit{user.sh} shell file is used, as the \textit{user.mk}, to
redefine shell functions.

With all these files a kaneton developer can parameterize the whole
kaneton project.

%
% interfaces
%

\section{Interfaces}

%
% make
%

\subsection{Make}

In this section we will detail the make interface that every machine
profile must implement.

The reader should look closer to the machine profiles already implemented.

The current machine reference profile is the \textit{linux} one.

\function{print}{(\argument{color},
                  \argument{text},
                  \argument{options})}
         {
	   This function display the \argument{text} in the specified
	   \argument{color}.

	   The option \textit{--no-newline} can be used not to output
	   the trailing newline.
	 }

\function{change-directory}{(\argument{directory},
                             \argument{options})}
         {
	   This function changes the current working directory.
	 }

\function{launch-shell}{(\argument{shell script},
                         \argument{arguments},
                         \argument{options})}
         {
	   This function launches a new shell script with its arguments.
	 }

\function{launch-python}{(\argument{python script},
                          \argument{arguments},
                          \argument{options})}
         {
	   This function launches a new python script with its arguments.
	 }

\function{launch-perl}{(\argument{shell script},
                        \argument{arguments},
                        \argument{options})}
         {
	   This function launches a new perl script with its arguments.
	 }

\function{launch-make}{(\argument{make file},
                        \argument{directory list},
                        \argument{arguments},
                        \argument{options})}
         {
	   This function launches the make command in each directory of
	   the list.
	 }

\function{launch}{(\argument{file},
                   \argument{...})}
         {
	   This function is a wrapper over the functions to launch new
	   shell, python, perl and make scripts.
	 }

\function{preprocess}{(\argument{preprocessed file},
                       \argument{c file},
                       \argument{options})}
         {
	   This function launches the C preprocessor.
	 }

\function{compile-c}{(\argument{object file},
                      \argument{c file},
                      \argument{options})}
         {
	   This function compile a C file generating an object file.
	 }

\function{lex-l}{(\argument{c file},
                  \argument{lex file},
                  \argument{options})}
         {
	   This function generates a C source file from a lex file.
	 }

\function{assemble-S}{(\argument{object file},
                       \argument{S file},
                       \argument{options})}
         {
	   This function assemble a S file.
	 }

\function{assemble-asm}{(\argument{object file},
                         \argument{asm file},
                         \argument{options})}
         {
	   This function assemble an asm file.

	   The option \textit{--output-format [format]} can be used to
	   specify the output format \textit{elf}, \textit{bin} etc..
	 }

\function{static-library}{(\argument{static library file name},
                           \argument{object files},
                           \argument{options})}
         {
	   This function builds a static library from object files.
	 }

\function{dynamic-library}{(\argument{dynamic library file name},
                            \argument{object files and/or libraries},
                            \argument{options})}
         {
	   This function builds a dynamic library from object files and/or
	   libraries.
	 }

\function{executable}{(\argument{executable file name},
                       \argument{object files and/or libraries},
                       \argument{options})}
         {
	   This function builds a executable from object files and/or
	   libraries.

	   The option \textit{--link-script [file]} specifies the
	   linker script file to use when building the executable object.

	   The option \textit{--no-standard-include} tells the function
	   not to use the operating system standard includes.

	   The option \textit{--no-standard-library} tells the function
	   not to use the operating system standard library.
	 }

\function{archive}{(\argument{archive file name},
                    \argument{object files},
                    \argument{options})}
         {
	   This function builds an archive of object files from multiple
	   object files.
	 }

\function{remove}{(\argument{files},
                   \argument{options})}
         {
	   This function removes the files in the list.
	 }

\function{purge}{()}
         {
	   This function just cleans the current working directory from
	   unecessary files.
	 }

\function{prototypes}{(\argument{files},
                       \argument{options})}
         {
	   This function will generate the prototypes in relation with
	   the files in the list.
	 }

\function{dependencies}{(\argument{files},
                         \argument{output},
                         \argument{options})}
         {
	   This function will generate the dependencies for the
	   \argument{files} and build a makefile dependency
	   \argument{output}.
	 }

\function{version}{(\argument{file})}
         {
	   This function generates a version file from the operating
	   system's informations: user, host, date etc..
	 }

\function{link}{(\argument{link},
                 \argument{file},
                 \argument{options})}
         {
	   This function creates a link \argument{link} for the file
	   \argument{file}.
	 }

\function{compile-tex}{(\argument{file},
                        \argument{options})}
         {
	   This function compiles the file \argument{file}.tex and
	   will generate a readable document.
	 }

\function{paper}{(\argument{file},
                  \argument{options})}
         {
	   This function builds a paper by calling the
	   \textbf{compile-tex()} function.
	 }

\function{lecture}{(\argument{file},
                    \argument{options})}
         {
	   This function builds a lecture document by calling the
	   \textbf{compile-tex()} function.
	 }

\function{subject}{(\argument{file},
                    \argument{options})}
         {
	   This function builds a subject by calling the
	   \textbf{compile-tex()} function.
	 }

\function{correction}{(\argument{file},
                       \argument{options})}
         {
	   This function builds a correction document by calling the
	   \textbf{compile-tex()} function.
	 }

\function{view}{(\argument{file},
                 \argument{options})}
         {
	   This function launches a viewer for the readable document
	   generated by the function \textbf{compile-tex()}.
	 }

\function{contents}{(\argument{file},
                     \argument{options})}
         {
	   This function returns the contents of the file \argument{file}.
	 }

%
% shell
%

\subsection{Shell}

In this section we will detail the shell interface that every machine
profile must implement.

The reader should look closer to the machine profiles already implemented.

\function{print}{(\argument{color},
                  \argument{text},
                  \argument{options})}
         {
	   This function outputs some text to the screen using a specific
	   color \textit{red}, \textit{green}, \textit{blue} etc..

	   The option \textit{--no-newline} can be used not to output
	   the trailing newline.

	   The option \textit{--flickering} tells the function to make
	   the text flicker.
	 }

\function{contents}{(\argument{file},
                     \argument{options})}
         {
	   This function returns the contents of the file \argument{file}.
	 }

\function{temporary}{(\argument{options})}
         {
	   This function creates a temporary file system object.

	   The option \textit{--file} tells the function to create a
	   temporary file.

	   The option \textit{--directory} tells the function to create
	   a temporary directory.
	 }

\function{working-directory}{(\argument{options})}
         {
	   This function returns the path of the current working directory.
	 }

\function{locate}{(\argument{name},
                   \argument{options})}
         {
	   This function tries to locate the program \argument{name}
	   on the system.
	 }

\function{substitute}{(\argument{pattern},
                       \argument{string},
                       \argument{options})}
         {
	   This function substitutes any string of the standard input
	   matching the pattern \argument{pattern} by the string
	   \argument{string}.

	   The option \textit{--everywhere} tells the function to
	   maximize the substitutions on the same line.
	 }

\function{display}{(\argument{message},
                    \argument{header})}
         {
	   This function displays the text \argument{message} in a
	   very specific manner based on the argument \argument{header}:
	   \textit{+}, \textit{!}, \textit{?} etc..
	 }

\function{wait-key}{(\argument{options})}
         {
	   This function just waits for the user to press a key.
	 }

\function{launch-shell}{(\argument{file},
                         \argument{arguments},
                         \argument{options})}
         {
	   This function launches a new shell script \argument{file} with
	   the arguments \argument{arguments}.
	 }

\function{launch-python}{(\argument{file},
                          \argument{arguments},
                          \argument{options})}
         {
	   This function launches the python file \argument{file} with
	   the arguments \argument{arguments}.
	 }

\function{launch-perl}{(\argument{file},
                        \argument{arguments},
                        \argument{options})}
         {
	   This function launches the perl file \argument{file} with
	   the arguments \argument{arguments}.
	 }

\function{launch-make}{(\argument{file},
                        \argument{directories},
                        \argument{arguments},
                        \argument{options})}
         {
	   This function launches the make file \argument{file} in each
	   directory of \argument{directories} with the arguments
	   \argument{arguments}.
	 }

\function{launch-program}{(\argument{file},
                           \argument{arguments},
                           \argument{options})}
         {
	   This function launches the program file \argument{file} with
	   the arguments \argument{arguments}.
	 }

\function{launch}{(\argument{file},
                   \argument{...})}
         {
	   This function is a wrapper over the previously described
	   launch functions.
	 }

\function{copy}{(\argument{source},
                 \argument{destination},
                 \argument{options})}
         {
	   This function copies the file \argument{source} to
	   \argument{destination}.
	 }

\function{link}{(\argument{source},
                 \argument{destination},
                 \argument{options})}
         {
	   This function builds a link between the file \argument{source}
	   and the file \argument{destination}.
	 }

\function{remove}{(\argument{file},
                   \argument{options})}
         {
	   This function removes the file \argument{file}.
	 }

\function{list}{(\argument{directory},
                 \argument{options})}
         {
	   This function lists the file system objects contains in the
	   directory \argument{directory}.

	   The option \textit{--file} tells the function to list the
	   files only.

	   The option \textit{--directory} tells the function to list
	   the directories only.
	 }

\function{change-directory}{(\argument{directory},
                             \argument{options})}
         {
	   This function changes the current working directory to
	   \argument{directory}.
	 }

\function{preprocess}{(\argument{file},
                       \argument{includes},
                       \argument{options})}
         {
	   This function performs a preprocessing task on the file
	   \argument{file}.

	   The argument \argument{includes} is used to provide additional
	   includes.

	   The option \textit{--no-markers} inhibits the generation
	   of linemarkers in the output.
	 }

\function{search}{(\argument{path},
                   \argument{pattern},
                   \argument{options})}
         {
	   This function looks for a file given the pattern
	   \argument{pattern} from the directory \argument{path}.

	   The option \textit{--file} is used to retrieve only files.

	   The option \textit{--directory} is used to retrieve only
	   directories.
	 }

\function{pack}{(\argument{directory},
                 \argument{file},
                 \argument{options})}
         {
	   This function makes an archive \argument{file} of the
	   directory \argument{directory}.
	 }

\function{unpack}{(\argument{directory},
                   \argument{file},
                   \argument{options})}
         {
	   This function extractes the archive \argument{file} into the
	   directory \argument{directory}.
	 }

\function{make-directory}{(\argument{directory},
                           \argument{options})}
         {
	   This function builds a new directory named
	   \argument{directory}.
	 }

\function{make-directory}{(\argument{directory},
                           \argument{options})}
         {
	   This function builds a new directory named
	   \argument{directory}.
	 }

\function{cut}{(\argument{head},
                \argument{tail})}
         {
	   This function tries to locate text between the two patterns
	   \argument{head} and \argument{tail} from the standard input.

	   The option \textit{--delete} tells the function to delete
	   this text and display the rest.

	   The option \textit{--print} tells the function to print
	   this text.
	 }

\function{load}{(\argument{source},
                 \argument{device},
                 \argument{path},
                 \argument{options})}
         {
	   This function loads the file \argument{source} on the
	   specificed device \argument{device}, more precisly
	   at the location \argument{path}.

	   The option \textit{--image} specifies that the argument
	   \argument{device} represents a device image.

	   The option \textit{--device} specifies that the argument
	   \argument{device} represents a device.
	 }

\function{stamp}{(\argument{format},
                  \argument{options})}
         {
	   This function returns a date following the given format
	   \argument{format}.
	 }

\function{stamp}{(\argument{format},
                  \argument{options})}
         {
	   This function returns a date following the given format
	   \argument{format}.
	 }

\function{record}{(\argument{log},
                   \argument{time},
                   \argument{options})}
         {
	   This function starts recording a session and outputs
	   the text into the file \argument{log} while the timings
	   are output in the file \argument{time}.
	 }

\function{play}{(\argument{log},
                 \argument{time},
                 \argument{options})}
         {
	   This function plays a previously recorded session where
	   the files \argument{log} and \argument{time} hold the
	   text and the timings.
	 }

%
% actions
%

\section{Actions}

Almost every different actions can be triggered from the kaneton
root directory.

These actions are listed below.

\command{make init}
        {
	  This action will initialize the kaneton development environment.

	  \example{\$ make init}
	}

\command{make clean}
	{
	  This action will clean the kaneton development environment.

	  \example{\$ make clean}
	}

\command{make kaneton}
	{
	  This action will build the kaneton microkernel.

	  \example{\$ make kaneton}
	}

\command{make clear}
	{
	  This action will removing every compiled files.

	  \example{\$ make clear}
	}

\command{make purge}
	{
	  This action will clean directories from unecessary files.

	  \example{\$ make purge}
	}

\command{make dep}
	{
	  This action will generate makefiles dependencies.

	  \example{\$ make dep}
	}

\command{make proto}
	{
	  This action will generate C prototypes.

	  \example{\$ make proto}
	}

\command{make check}
	{
	  This action will be used to run the test suite used to validate
	  the kaneton whole behaviour.

	  \example{\$ make check}
	}

\command{make cheat}
	{
	  This action will launch the cheat tests on students kaneton
	  implementations.

	  \example{\$ make cheat}

	  \example{\$ make cheat-2006.k3}
	}

\command{make build}
	{
	  This action will build the boot device.

	  \example{\$ make build}
	}

\command{make install}
	{
	  This action will install the kaneton stuff on the boot device.

	  \example{\$ make install}
	}

\command{make export}
	{
	  This action will build a kaneton package.

	  \example{\$ make export}

	  \example{\$ make export-k5}
	}

\command{make view}
	{
	  This action will generate and display a kaneton document.

	  \example{\$ make view}

	  \example{\$ make view-develop}

	  \example{\$ make view-papers::kaneton}
	}

\command{make record}
	{
	  This action will record a real-time session.

	  \example{\$ make record}

	  \example{\$ make record-sessions::test.ts}
	}

\command{make play}
	{
	  This action will play a recorded session.

	  \example{\$ make play}

	  \example{\$ make play-sessions::prototypes.ts}
	}

\command{make info}
	{
	  This action will display general kaneton information.

	  \example{\$ make info}
	}

%\include{coding}
%%
% ---------- header -----------------------------------------------------------
%
% project       kaneton
%
% license       kaneton
%
% file          /home/mycure/kaneton/view/book/core/licenses.tex
%
% created       julien quintard   [tue may 12 18:35:31 2009]
% updated       julien quintard   [tue may 12 18:35:33 2009]
%

%
% ---------- licenses ---------------------------------------------------------
%

\chapter{Licenses}
\label{chapter:licenses}

In this chapter we will details the kaneton-related licenses.

\newpage

%
% ---------- text -------------------------------------------------------------
%

The kaneton project might be considered as an open-project since source
code is provided.

Nevertheless this is not the case as this project is used as material for
operating system courses.

Therefore, people implementing the kaneton microkernel should not
make their source code available. To avoid problems, especially students
cheating, kaneton people decided to use a kaneton-specific license
forbidding source code distribution.

The \name{kaneton license} is based on a more generic license, the
\name{pedagogical licence}.

The next sections will contain these licenses' descriptions.

%
% pedagogical license
%

\section{Pedagogical License}

\textbf{TERMS AND CONDITIONS FOR COPYING, DISTRIBUTION AND MODIFICATION}

\begin{enumerate}
  \item
    This License applies to any program or other work which contains a
    notice placed by the copyright holder saying it may be distributed
    under the terms of this License. The "Program", below, refers
    to any such program or work, and a "work based on the Program"
    means either the Program or any derivative work under copyright law:
    that is to say, a work containing the Program or a portion of it,
    either verbatim or with modifications and/or translated into another
    language. (Hereinafter, translation is included without limitation in
    the term "modification".) Each licensee is addressed as "you".
  \item
    You must not copy or distribute copies of the Program's source
    code, object code or executable form without explicit authorization
    from the maintainers.

    If you have this authorization, you must conspicuously and
    appropriately publish on each copy an appropriate copyright notice
    and disclaimer of warranty; keep intact all the notices that refer
    to this License and to the absence of any warranty; and give any
    other recipients of the Program a copy of this License along with
    the Program.
  \item
    You may modify your copy or copies of the Program or any portion of
    it, thus forming a work based on the Program, provided that you
    meet all of these conditions :

    \begin{enumerate}
      \item
	You must not publish your work without explicit authorization
	from the maintainers.
      \item
	You must send to the maintainers any work that in whole or in
	part contains or is derived from the Program or any part thereof.
      \item
	You must cause any work that you send, that in whole or in
	part contains or is derived from the Program or any part thereof,
	to be licensed as a whole under the terms of this License.
      \item
	You must cause the modified files to carry prominent notices
	stating that you changed the files and the date of any change.
      \item
	If the modified program normally reads commands interactively
	when run, you must cause it, when started running for such
	interactive use in the most ordinary way, to print or display an
	announcement including an appropriate copyright notice and a
	notice that there is no warranty (or else, saying that you
	provide a warranty) and telling the user how to view a copy
	of this License. (Exception: if the Program itself is interactive
	but does not normally print such an announcement, your work based
	on the Program is not required to print an announcement.)
    \end{enumerate}
  \item
    Access to the Program's source is granted if either:

    \begin{enumerate}
      \item
	You want to make the Program evolve
      \item
	You have pedagogical goals
    \end{enumerate}

    These requirements apply to the modified work as a whole. If
    identifiable sections of that work are not derived from the Program,
    and can be reasonably considered independent and separate works in
    themselves, then this License, and its terms, do not apply to those
    sections when you distribute them as separate works. But when you
    distribute the same sections as part of a whole which is a work based
    on the Program, the distribution of the whole must be on the terms of
    this License, whose permissions for other licensees extend to the
    entire whole, and thus to each and every part regardless of who wrote
    it.

    Thus, it is not the intent of this section to claim rights or contest
    your rights to work written entirely by you; rather, the intent is to
    exercise the right to control the distribution of derivative or
    collective works based on the Program.

    In addition, mere aggregation of another work not based on the Program
    with the Program (or with a work based on the Program) on a volume of
    a storage or distribution medium does not bring the other work under
    the scope of this License.
  \item
    You may not copy, modify, sublicense, or distribute the Program
    except as expressly provided under this License. Any attempt
    otherwise to copy, modify, sublicense or distribute the Program is
    void, and will automatically terminate your rights under this
    License. However, parties who have received copies, or rights,
    from you under this License will not have their licenses
    terminated so long as such parties remain in full compliance.
  \item
    You are not required to accept this License, since you have not
    signed it. However, nothing else grants you permission to modify
    or distribute the Program or its derivative works. These actions
    are prohibited by law if you do not accept this
    License. Therefore, by modifying or distributing the Program (or
    any work based on the Program), you indicate your acceptance of
    this License to do so, and all its terms and conditions for
    copying, distributing or modifying the Program or works based on
    it.
  \item
    Each time you redistribute the Program (or any work based on the
    Program), the recipient automatically receives a license from the
    original licensor to copy, distribute or modify the Program subject
    to these terms and conditions. You may not impose any further
    restrictions on the recipients' exercise of the rights granted
    herein. You are not responsible for enforcing compliance by third
    parties to this License.
  \item
    If the distribution and/or use of the Program is restricted in
    certain countries either by patents or by copyrighted interfaces,
    the original copyright holder who places the Program under this
    License may add an explicit geographical distribution limitation
    excluding those countries, so that distribution is permitted only
    in or among countries not thus excluded. In such case, this License
    incorporates the limitation as if written in the body of this
    License.
  \item
    We may publish revised and/or new versions of this License from
    time to time. It may evolve considering new contributors
    needs. Contact us if you have any request. Such new versions will
    be similar in spirit to the present version, but may differ in
    detail to address new problems or concerns.
  \item
    If you wish to incorporate parts of the Program into other free
    programs whose distribution conditions are different, write to the
    author to ask for permission, we sometimes make exceptions for this.
  \item
    THERE IS NO WARRANTY FOR THE PROGRAM, TO THE EXTENT PERMITTED BY
    APPLICABLE LAW. EXCEPT WHEN OTHERWISE STATED IN WRITING THE
    COPYRIGHT HOLDERS AND/OR OTHER PARTIES PROVIDE THE PROGRAM "AS IS"
    WITHOUT WARRANTY OF ANY KIND, EITHER EXPRESSED OR IMPLIED,
    INCLUDING, BUT NOT LIMITED TO, THE IMPLIED WARRANTIES OF
    MERCHANTABILITY AND FITNESS FOR A PARTICULAR PURPOSE. THE ENTIRE
    RISK AS TO THE QUALITY AND PERFORMANCE OF THE PROGRAM IS WITH
    YOU. SHOULD THE PROGRAM PROVE DEFECTIVE, YOU ASSUME THE COST OF
    ALL NECESSARY SERVICING, REPAIR OR CORRECTION.
  \item
    IN NO EVENT UNLESS REQUIRED BY APPLICABLE LAW OR AGREED TO IN
    WRITING WILL ANY COPYRIGHT HOLDER, OR ANY OTHER PARTY WHO MAY
    MODIFY AND/OR REDISTRIBUTE THE PROGRAM AS PERMITTED ABOVE, BE
    LIABLE TO YOU FOR DAMAGES, INCLUDING ANY GENERAL, SPECIAL,
    INCIDENTAL OR CONSEQUENTIAL DAMAGES ARISING OUT OF THE USE OR
    INABILITY TO USE THE PROGRAM (INCLUDING BUT NOT LIMITED TO LOSS OF
    DATA OR DATA BEING RENDERED INACCURATE OR LOSSES SUSTAINED BY YOU
    OR THIRD PARTIES OR A FAILURE OF THE PROGRAM TO OPERATE WITH ANY
    OTHER PROGRAMS), EVEN IF SUCH HOLDER OR OTHER PARTY HAS BEEN
    ADVISED OF THE POSSIBILITY OF SUCH DAMAGES.
\end{enumerate}

%
% kaneton license
%

\section{kaneton License}

\textbf{TERMS AND CONDITIONS FOR COPYING, DISTRIBUTION AND MODIFICATION}

\begin{enumerate}
  \item
    This licence is nothing more than a link to the pedagogical licence.
  \item
    Any program under the kaneton licence is in fact under the terms and
    conditions of the pedagogical licence.
\end{enumerate}

%%%
%% licence       kaneton licence
%%
%% project       kaneton
%%
%% file          /home/mycure/kaneton/view/papers/kaneton/bibliography.tex
%%
%% created       julien quintard   [mon may  8 18:35:35 2006]
%% updated       julien quintard   [mon may  8 20:38:56 2006]
%%

%
% bibliograpy
%

\chapter{Bibliography}

This chapter contains the bibliography.

%
% text
%

\begin{thebibliography}{0}
  \bibitem{AST-SCO}
    \textbf{Structured Computer Organization};
    by
    \textit{Andrew S. Tanenbaum}
  \bibitem{AST-CN}
    \textbf{Computer Networks};
    by
    \textit{Andrew S. Tanenbaum}
  \bibitem{AST-OSDI}
    \textbf{Operating Systems: Design and Implementation};
    by
    \textit{Andrew S. Tanenbaum, Albert S Woodhull}
  \bibitem{AST-MOS}
    \textbf{Modern Operating Systems};
    by
    \textit{Andrew S. Tanenbaum}
  \bibitem{AST-DOS}
    \textbf{Distributed Operating Systems};
    by
    \textit{Andrew S. Tanenbaum}
  \bibitem{AST-DSPP}
    \textbf{Distributed Systems: Principles and Paradigms};
    by
    \textit{Andrew S. Tanenbaum, Maarten van Steen}
  \bibitem{NAL-DA}
    \textbf{Distributed Algorithms};
    by
    \textit{Nancy A. Lynch}
\end{thebibliography}


\end{document}

%
% ---------- intro ------------------------------------------------------------
%

kaneton est un projet pedagogique -> code clair, documente, cours etc..

un point tres important est la clarete du code. cela implique une rigeur
exceptionnel de la part des developpeurs.

le projet est compose du code du kernel, d'un certain nombre de scripts
et makefiles pour faciliter bcp d'operations et de documentations.

ce document decrit donc en profondeurs les regles qu'un contributeur doit
absolument suivre mais revient egalement sur la philisophie d'un travail
communautaire qui implique un comportement exemplaire de la part de
chaque developpeurs.

%
% ---------- community --------------------------------------------------------
%

regles communautaire.

-- comportement

mieux plutot que meilleur

savoir participer en faisant des trucs moins cool que d'autres car ca
doit etre fait

...

le projet dispose de certains outils pour effectuer le minimum vital:
outils de comm interne, dev communautaire, outils de comm ext, gestion
de taches.

%
% ---------- tools ------------------------------------------------------------
%

regles concernant les outils.

-- communication

parceque c'est chiant d'avoir plusieurs outils de comm mail, news, forum
et que tout le monde dispose d'un mail, on utilise une ml.

comportement tres important sur une ml: eviter les trolls, repondre
en suvant la netiquette, participer en donnant son avis etc..

pour la communication temps-reel, un serveur irc pourrait etre mis en
place mais ce n'est pas le cas actuellemeent. si vous en ressentez le
desir, proposez l'idee sur la ml. dans tous les cas, cela ne sera pas
obligatoire.

de maniere generale, des outils libres, pre-fournis et simple d'utilisation
seront privlegies. ainsi google fournit un certains nombre d'outils que
kaneton utilise: gmail opur les mails si vous le desirez, googlegroup pour
la ml, project hosting pour le projet ..

-- repository

le repository est l'outil le plus important du projet puisque tout y
est stocke, le code bien sur mais egalement les documentations, les
projets des etudiants au fil des annees ...

le repository actuel contient toute l'histoire du projet kaneton depuis
le premier jour jusqu'a aujourdhui. nous utilisons subversion car a l'epoque
c'etait de loin le meilleur outil, son concurrent CVS etant deja trop
limite.

durant les premieres annees aucune normalisation de l'utilisation du
repository ne fut instauree. maintenant, les commits doivent suivre une
procedure praticuliere qui implique un descriptif normalise.

de plus l'organisation du repository est tres tres important et chaque ajout
doit etre fait intelligemment afin d'eviter que quelqu'un soit oblige de
repasser derriere pour reparer ces erreurs.

-- wiki

le wiki contient trois parties: une partie publique et deux wikis: un pour
les etudiants et un pour les developpeurs.

le wiki kaneton.org doit etre utilise en priorite du wiki associe avec
le systeme de tickets etc..

participer au wiki est tres important puisqu'il permet de mettre en avant
l'evolution du projet.

neanmoins suivant le nombre de developpeurs impliques il est normal que
certains soient moins actifs que d'autres mais neanmoins participer a
la documentation et donc au wiki est aussi important que le code lui-meme.

-- tickets

le systeme de tickets/bugs est egalement tres important. chaque ticket se
voit affecte une priorite et il est important de comprendre que pour le
bien global du projet, un developpeur ne peut se contenter de faire ce
qui lui plait, il se doit de contribuer egalement a la resolution de problemes.

encore une fois les tickets doivent suivre une norme.

%
% ---------- languages --------------------------------------------------------
%

pour chaque language explique les regles: header, commentaire etc..

-- c

-- make

-- python

-- latex

figures in .fig

-- asm

-- shell
